\scanpage{78}
\margin{\textcircled{2}}
\uline{Exemples \inserted{de corps} d’anneaux} (Cas de corps \deleted{\unsure{de}} de car $\neq2$)

\diagram{\begin{array}{llll}
\text{a) \uline{corps alg. clos} }k&F(k)\simeq\mathbf Z&[\text{\unsure{i.e.}}\ F_1(k)=0\ ]\\[3pt]
\text{b) \uline{Corps fini} }K&F(k)\simeq\mathbf Z+\mathbf Z/2\mathbf Z&[\text{\unsure{i.e.}}\ F_2^{\deleted{3}}(k)=0\ ]\\[3pt]
\text{\unsure{b|c}) \uline{Corps} }\uline{p\text{-adiques}}\ K&F(K)\simeq\mathbf Z\ \begin{array}{l}\text{étendu par un}\\\text{groupe fini}\\\text{annulé par $4$}\end{array}&[\text{\unsure{i.e.}}\ F_3(k)\unsure{=0}\ ]\\[3pt]
\text{d) \uline{Corps} }\mathbf R&\begin{array}{l}F_1(k)\simeq\mathbf Z\\\text{par l’index}\\F_2(k)=2\mathbf Z\\F_3(k)=4\mathbf Z\end{array}&\deleted{\unsure{F_2}}
\end{array}}

e) Corps valué complet à corps résiduel alg clos \deleted{\unsure{discrètement}}
\[
F_{\unsure{1|k}}(K)\simeq\mathbf Z+\underbrace{K^*/k^{*2}}_{\mathbf Z/2\mathbf Z}\qquad[\text{On a }F_2(k)=0]
\]
\margin{si car. résiduelle est $\neq2$, \unsure{quasi par définition...}}

\margin{\textcircled{3}}
\deleted{\unsure{3}} \uline{Corps de nombres}

\uline{Th. de Hasse}
\diagram{\begin{tikzcd}[column sep=large]
F(K)\arrow[r]&\sum_pF(K_p)&\text{\uuline{injectif}}\\
F_1(K)/F_2(K)\arrow[r]&\sum_pK_p^*/K_p^{*2}&\begin{array}{l}\text{injectif [grâce à Hasse}\\\text{\unsure{égalemt}] }\end{array}\\
\underline{F_2(K)/F_3(K)}\arrow[r]&\sum_p\mathbf Z/2\mathbf Z&\begin{array}{l}\text{injectif [car }\operatorname{Br}(K)\\\subset\sum_p\operatorname{Br}(K_p)]\end{array}
\end{tikzcd}}
\inserted{sous-groupe des \unsure{formes d’invariant de norme} unité}
\[
F_3(K)\subset\sum_{\substack{p\text{ infini}\\\text{réel}}}F_3(K_p)\subset\mathbf Z^r
\]
\margin{$r$ nb de places réelles}

Le sous-groupe de torsion de $F(K)$ est le noyau de
\diagram{\begin{tikzcd}
F_1(K)\arrow[r,"\text{index}","\substack{\text{aux places}\\\text{réelles}}"']&\mathbf Z^r
\end{tikzcd}}
, il est annulé par $4$
\margin{77}
